% Part: first-order-logic
% Chapter: introduction
% Section: semantic-notions

\documentclass[../../../include/open-logic-section]{subfiles}

\begin{document}

\olfileid{fol}{int}{sem}

\olsection{معنايي مفاهيم}

پورته مو يادونه وکړه چې کله ګورو ايا $\Sat{M}{!A}[s]$ صدق کوي، د
اسانتيا دپاره~$s$ ته پرېږدو چې ټولو !!{variable}s ته قيمتونه
وګوماري، خو يوازې هغه قيمتونه کارېږي چې په~$!A$ کښې
!!{variable}s ته ئې ګوماري۔ په حقيقت کښې، يوازې په~$!A$ کښې د
\emph{ازادو} متغيرونو قيمتونه اهميت لري۔ البته ځکه چې موږ احتياط
کوو، دا حقيقت به ثابتوو۔ ځکه چې !!{sentence}s ازاد متغيرونه نۀ
لري، نو په !!a{structure} کښې د هغوئ د صدق دپاره~$s$ هېڅ اهميت نۀ
لري۔ نو کله چې~$!A$ !!a{sentence} وي، $\Sat{M}{!A}$ داسې تعريفولے
شو چې معنٰی ئې ``د ټولو~$s$ دپاره $\Sat{M}{!A}[s]$'' ده؛ او تصادفاً
دا هله او يوازې هله رښتيا دي چې لږ تر لږه د يوه~$s$ دپاره
$\Sat{M}{!A}[s]$ وي۔ د !!{formula}s د صدق د کار وړ تعريف دپاره د
!!{variable} ګومارنې ورپېژندل پکار دي، خو د !!{sentence}s صدق د
!!{variable} له ګومارنو څخه خپلواک دے۔

کله چې د ``$\Sat{M}{!A}$'' تعريف ولرو، د لومړۍ درجې منطق د
!!{sentence}s په باب پوهېږو چې ``حالت'' او ``پکښې رښتونې'' څه معنٰی
لري۔ د !!{sentence}s دپاره د~$\Sat{M}{!A}$ د تعريف پر بنسټ بيا د
اعتبار، استلزام او صدق وړتيا بنسټيز معنايي مفاهيم تعريفولے شو۔ يوه
جمله، $\Entails !A$، هله معتبره ده چې هره !!{structure} ئې پوره
کړي۔ د !!{sentence}s د يوه سټ له خوا، $\Gamma \Entails !A$،
هله لازمېږي چې هره هغه !!{structure} چې په~$\Gamma$ کښې ټولې
!!{sentence}s پوره کوي،~$!A$ هم پوره کړي۔ او د !!{sentence}s يو سټ
هله د صدق وړ دے چې کومه !!{structure} پکښې ټولې !!{sentence}s په
يو وخت پوره کړي۔

ځکه چې !!{formula}s په استقرا تعريف شوي، او صدق بيا د !!{formula}s
پر جوړښت په استقرا تعريفېږي، د خپلې معنٰی پوهنې د خاصيتونو د ثبوت
او د تعريف شوو معنايي مفاهيمو د تړلو دپاره استقرا کارولے شو۔ ځينې
دغه خاصيتونه به راټول او ثابت کړو؛ يوه وجه دا ده چې هر يو پخپله په
زړه پورې دے، خو اصلي وجه دا ده چې وروسته د معنٰی پوهنې او د
!!{derivation} د نظامونو د اړيکې په څېړنه کښې ډېر راپکارېږي۔ د دې
کار دپاره به ځينې داسې نحوي مفاهيم او عمليات هم په دقيق ډول (يعنې
په استقرا) تعريفول پکار وي چې لا مو نۀ دي ياد کړي۔

\end{document}
